Esta seção retoma o escopo dos cinco épicos definidos na Seção~\ref{sec:requisitos} e descreve o que deve ser desenvolvido em cada camada do sistema. As funcionalidades correspondem às histórias de usuário especificadas no Capítulo~\ref{cap:metodologia}. As tabelas do backend relacionam essas histórias às operações atualmente documentadas no Swagger da API Seed. O mapeamento é feito pela finalidade descrita para cada operação.

\subsection{Épico 1 -- Gerenciamento de Registro do Usuário}

Este épico reúne as funcionalidades de cadastro e manutenção dos dados de usuários e funcionários, previstas nas histórias US1 a US4.

\subsubsection{\textit{Frontend}}

Devem ser disponibilizadas telas responsivas para que usuários e funcionários realizem seu cadastro e editem seus dados. Também deve ser oferecida uma interação para solicitar a exclusão dos dados. Os formulários devem permitir informar os dados necessários a cada operação e apresentar o resultado da solicitação ao usuário.

\subsubsection{\textit{Backend}}

As operações abaixo foram localizadas no Swagger e associadas às histórias conforme sua finalidade. Com exceção das rotas de cadastro de estudante, cadastro de funcionário e \textit{login}, as demais rotas da API exigem um token JWT. O \textit{login} é realizado por meio de \textit{POST} em \texttt{/auth/login} e, quando bem-sucedido, retorna o token que será enviado nas requisições seguintes para autenticar e autorizar o usuário. O token carrega três informações:

\begin{itemize}
    \item \texttt{UUID}: identificador único que permite localizar os dados pessoais do usuário;
    \item \texttt{ROLE}: define o papel e as permissões dele;
    \item \texttt{EXPIRATION}: data de expiração do token.
\end{itemize}

Quando a operação envolve dados pessoais, o UUID indica quem é o titular e restringe consultas e alterações aos dados do próprio usuário, de modo que ninguém consiga mexer nas informações de outra pessoa. A cada requisição, o \textit{backend} confere se o token ainda é válido e verifica se o UUID tem a permissão que aquela operação exige, no banco de dados, sendo assim, quem não tem a permissão registrada não consegue entrar no painel administrativo de uma IES. Já nas consultas e alterações de dados de uma IES, o \textit{backend} usa o UUID para descobrir a qual instituição o usuário está vinculado e limita a operação a ela, evitando que alguém acesse ou altere dados de outras IES.

A Tabela~\ref{tab:endpoints-autenticacao} relaciona as histórias de usuário com as operações da API.

\renewcommand{\arraystretch}{1.5}
\setlength{\LTleft}{\fill}
\setlength{\LTright}{\fill}
\begin{longtable}{|>{\centering\arraybackslash}p{1.2cm}|>{\centering\arraybackslash}p{1.7cm}|>{\raggedright\arraybackslash}p{5.7cm}|>{\raggedright\arraybackslash}p{5.2cm}|}
\caption{\textit{Endpoints} de autenticação.}\label{tab:endpoints-autenticacao}\\
\hline
\textbf{US} & \textbf{Método} & \textbf{Rota} & \textbf{Operação} \\
\hline
\endfirsthead
\multicolumn{4}{c}{\tablename\ \thetable{} -- continuação}\\
\hline
\textbf{US} & \textbf{Método} & \textbf{Rota} & \textbf{Operação} \\
\hline
\endhead
\hline
\endfoot
\hline
\multicolumn{4}{c}{\small Fonte: autoria própria.}\\
\endlastfoot
US1 & \textit{POST} & \texttt{/auth/\allowbreak cadastro/\allowbreak estudante} & Cadastrar estudante. \\
\hline
US2 & \textit{POST} & \texttt{/auth/\allowbreak cadastro/\allowbreak funcionario} & Cadastrar funcionário. \\
\hline
-- & \textit{POST} & \texttt{/auth/login} & Autenticar usuário e retornar o token JWT. \\
\hline
US3 & \textit{GET} & \texttt{/estudante} & Consultar os próprios dados de estudante. \\
\hline
US3 & \textit{PATCH} & \texttt{/estudante} & Editar os próprios dados de estudante. \\
\hline
US3 & \textit{GET} & \texttt{/funcionario} & Consultar os próprios dados de funcionário. \\
\hline
US3 & \textit{PATCH} & \texttt{/funcionario} & Editar os próprios dados de funcionário. \\
\hline
US4 & \textit{DELETE} & \texttt{/estudante} & Excluir a própria conta de estudante. \\
\hline
US4 & \textit{DELETE} & \texttt{/funcionario} & Excluir a própria conta de funcionário. \\
\hline
\end{longtable}

\subsection{Épico 2 -- Gerenciamento de Formulários}

Este épico abrange o registro diário das avaliações de bem-estar, a consulta das respostas, sua apresentação em estatísticas e o gerenciamento de formulários e anotações, conforme as histórias US5 a US11.

\subsubsection{\textit{Frontend}}

Devem ser criadas telas responsivas para o usuário registrar sua avaliação diária em cada seção de bem-estar e consultar formulários já respondidos. O funcionário deve poder visualizar os dados registrados pelos usuários e criar ou editar formulários personalizados. A tela de preenchimento deve permitir que o usuário inclua uma anotação livre junto com suas respostas. Como comportamento previsto para a PWA, o formulário diário poderá ser preenchido sem conexão, mantendo as respostas pendentes localmente para envio à API após o restabelecimento da conexão.

\subsubsection{\textit{Backend}}

As rotas documentadas para formulários base e personalizados estão relacionadas na Tabela~\ref{tab:endpoints-formularios}. As histórias US8 e US11 não possuem rotas próprias, pois a análise estatística será realizada pelo \textit{Power BI}, e as anotações serão enviadas junto com as respostas durante o preenchimento do formulário prevista na US5.

\renewcommand{\arraystretch}{1.5}
\setlength{\LTleft}{\fill}
\setlength{\LTright}{\fill}
\begin{longtable}{|>{\centering\arraybackslash}p{1.2cm}|>{\centering\arraybackslash}p{1.7cm}|>{\raggedright\arraybackslash}p{5.7cm}|>{\raggedright\arraybackslash}p{5.2cm}|}
\caption{\textit{Endpoints} de formulários.}\label{tab:endpoints-formularios}\\
\hline
\textbf{US} & \textbf{Método} & \textbf{Rota} & \textbf{Operação} \\
\hline
\endfirsthead
\multicolumn{4}{c}{\tablename\ \thetable{} -- continuação}\\
\hline
\textbf{US} & \textbf{Método} & \textbf{Rota} & \textbf{Operação} \\
\hline
\endhead
\hline
\endfoot
\hline
\multicolumn{4}{c}{\small Fonte: autoria própria.}\\
\endlastfoot
US5 & \textit{POST} & \texttt{/api/\allowbreak estudante/\allowbreak formularios-base/\allowbreak respostas} & Preencher o formulário base. \\
\hline
US6 & \textit{GET} & \texttt{/api/\allowbreak funcionario/\allowbreak formularios-base/\allowbreak respostas} & Listar respostas dos estudantes da IES do funcionário. \\
\hline
US6 & \textit{GET} & \texttt{/api/\allowbreak funcionario/\allowbreak formularios-personalizados/\allowbreak \{id\}/\allowbreak respostas} & Listar respostas dos estudantes para um formulário personalizado. \\
\hline
US6 & \textit{GET} & \texttt{/api/\allowbreak funcionario/\allowbreak formularios-base/\allowbreak respostas/\allowbreak anotacoes} & Consultar as anotações das respostas dos estudantes da IES. \\
\hline
US7 & \textit{GET} & \texttt{/api/\allowbreak estudante/\allowbreak formularios-base/\allowbreak respostas} & Listar as próprias respostas do formulário base. \\
\hline
US7 & \textit{GET} & \texttt{/api/\allowbreak estudante/\allowbreak formularios-base/\allowbreak respostas/\allowbreak \{id\}} & Consultar uma resposta própria do formulário base. \\
\hline
US7 & \textit{GET} & \texttt{/api/\allowbreak estudante/\allowbreak formularios-personalizados/\allowbreak \{id\}/\allowbreak respostas} & Listar as próprias tentativas de resposta de um formulário personalizado. \\
\hline
US9 & \textit{POST} & \texttt{/api/\allowbreak funcionario/\allowbreak formularios-personalizados} & Criar formulário personalizado para a IES do funcionário. \\
\hline
US9 & \textit{GET} & \texttt{/api/\allowbreak funcionario/\allowbreak perguntas-\allowbreak formularios-\allowbreak personalizados} & Pesquisar perguntas reutilizáveis para os formulários. \\
\hline
US10 & \textit{GET} & \texttt{/api/\allowbreak funcionario/\allowbreak formularios-personalizados/\allowbreak \{id\}} & Consultar formulário e perguntas. \\
\hline
US10 & \textit{PUT} & \texttt{/api/\allowbreak funcionario/\allowbreak formularios-personalizados/\allowbreak \{id\}} & Editar formulário personalizado. \\
\hline
\end{longtable}


\subsection{Épico 3 -- Gerenciamento de Comunicação}

Este épico reúne as comunicações previstas nas histórias US12 a US15, como lembretes para preenchimento dos formulários, confirmação de troca de senha, avisos sobre o \textit{status} do sistema e avisos de eventos.

\subsubsection{\textit{Frontend}}

Devem ser contempladas as interações necessárias para que o usuário solicite a troca de senha e receba o retorno da operação. Os lembretes de formulários, os avisos sobre o \textit{status} do sistema e as comunicações sobre eventos são enviados por e-mail pelo serviço de notificação; portanto, não dependem de telas específicas para disparar cada mensagem.

\subsubsection{\textit{Backend}}

%O serviço de notificação por e-mail deve ser acionado por eventos internos da aplicação, sem criar rotas específicas para cada tipo de aviso. Para a US12, o disparo ocorre na execução diária programada quando o formulário do dia ainda estiver pendente. Para a US13, a confirmação é enviada após a troca de senha ser concluída com sucesso. Para a US14, uma alteração de \textit{status} que deva ser comunicada aciona o envio do aviso correspondente. Para a US15, um evento do calendário que deva gerar comunicação aciona o envio aos usuários destinatários. A periodicidade e as condições exatas de envio devem seguir as regras definidas para cada funcionalidade.

O serviço de notificação por e-mail deve ser acionado por eventos internos da aplicação, sem criar rotas específicas para cada tipo de aviso. O disparo ocorre da seguinte forma:

\begin{itemize}
    \item \textbf{US12:} na execução diária programada, quando o formulário do dia ainda estiver pendente;
    \item \textbf{US13:} após a conclusão bem-sucedida da troca de senha, com o envio da confirmação;
    \item \textbf{US14:} quando houver uma alteração de \textit{status} que deva ser comunicada, o que aciona o envio do aviso correspondente;
    \item \textbf{US15:} quando um evento do calendário deva gerar comunicação, o que aciona o envio aos usuários destinatários.
\end{itemize}

%A periodicidade e as condições exatas de envio devem seguir as regras definidas para cada funcionalidade.

As únicas rotas da API relacionadas a esse fluxo são as de alteração de senha. Elas iniciam a operação solicitada pelo usuário; já o envio do e-mail de confirmação é uma ação interna, disparada após a conclusão bem-sucedida da troca, sem uma rota própria de notificação. Assim, a Tabela~\ref{tab:gatilhos-comunicacao} apresenta as rotas desenvolvidas para o Épico 3.

\renewcommand{\arraystretch}{1.5}
\setlength{\LTleft}{\fill}
\setlength{\LTright}{\fill}
\begin{longtable}{|>{\centering\arraybackslash}p{1.2cm}|>{\raggedright\arraybackslash}p{5.4cm}|>{\raggedright\arraybackslash}p{6.5cm}|}
\caption{Gatilhos de comunicação por e-mail.}\label{tab:gatilhos-comunicacao}\\
\hline
\textbf{US} & \textbf{Gatilho} & \textbf{Ação do serviço de e-mail} \\
\hline
\endfirsthead
\multicolumn{3}{c}{\tablename\ \thetable{} -- continuação}\\
\hline
\textbf{US} & \textbf{Gatilho} & \textbf{Ação do serviço de e-mail} \\
\hline
\endhead
\hline
\endfoot
\hline
\multicolumn{3}{c}{\small Fonte: autoria própria.}\\
\endlastfoot
US12 & Execução diária programada, quando o formulário do dia estiver pendente. & Enviar ao usuário um lembrete por e-mail para preencher os formulários diários. \\
\hline
US13 & Conclusão bem-sucedida de \textit{PATCH} \texttt{/estudante/\allowbreak senha} ou \textit{PATCH} \texttt{/funcionario/\allowbreak senha}. & Enviar ao usuário a confirmação da troca de senha por e-mail. \\
\hline
US14 & Alteração do \textit{status} do sistema que deva ser comunicada. & Enviar por e-mail um aviso com a informação atualizada sobre o sistema. \\
\hline
US15 & Evento do calendário que, conforme as regras da aplicação, deva gerar uma comunicação. & Enviar por e-mail um aviso sobre o evento aos usuários destinatários. \\
\hline
\end{longtable}

\subsection{Épico 4 -- Gerenciamento de Eventos}

Este épico compreende o gerenciamento de eventos do calendário pelo funcionário e a visualização desses eventos pelo usuário, conforme as histórias US17 a US20.

\subsubsection{\textit{Frontend}}

Devem ser disponibilizadas telas responsivas para que o funcionário registre, edite e exclua eventos do calendário. Para o usuário, deve ser apresentada uma forma de consultar os eventos disponíveis no calendário do sistema.

\subsubsection{\textit{Backend}}

As operações documentadas para criação, edição, exclusão e consulta de eventos são relacionadas às histórias correspondentes na Tabela~\ref{tab:endpoints-eventos}.

\renewcommand{\arraystretch}{1.5}
\setlength{\LTleft}{\fill}
\setlength{\LTright}{\fill}
\begin{longtable}{|>{\centering\arraybackslash}p{1.2cm}|>{\centering\arraybackslash}p{1.7cm}|>{\raggedright\arraybackslash}p{5.7cm}|>{\raggedright\arraybackslash}p{5.2cm}|}
\caption{\textit{Endpoints} de eventos.}\label{tab:endpoints-eventos}\\
\hline
\textbf{US} & \textbf{Método} & \textbf{Rota} & \textbf{Operação} \\
\hline
\endfirsthead
\multicolumn{4}{c}{\tablename\ \thetable{} -- continuação}\\
\hline
\textbf{US} & \textbf{Método} & \textbf{Rota} & \textbf{Operação} \\
\hline
\endhead
\hline
\endfoot
\hline
\multicolumn{4}{c}{\small Fonte: autoria própria.}\\
\endlastfoot
US17 & \textit{POST} & \texttt{/eventos} & Criar evento. \\
\hline
US18 & \textit{PUT} & \texttt{/eventos/\allowbreak \{id\}} & Atualizar evento. \\
\hline
US19 & \textit{DELETE} & \texttt{/eventos/\allowbreak \{id\}} & Excluir evento. \\
\hline
US20 & \textit{GET} & \texttt{/eventos} & Listar eventos para consulta. \\
\hline
US20 & \textit{GET} & \texttt{/eventos/\allowbreak \{id\}} & Consultar evento por identificador. \\
\hline
\end{longtable}

\subsection{Épico 5 -- Gerenciamento do \textit{Feed}}

Este épico prevê a publicação e a consulta de conteúdo no \textit{feed} de cada seção de bem-estar, conforme as histórias US21 a US24.

\subsubsection{\textit{Frontend}}

Devem ser disponibilizadas telas responsivas para que o funcionário registre, edite e exclua publicações vinculadas às seções de bem-estar. O usuário deve poder visualizar o \textit{feed} correspondente a cada seção.

\subsubsection{\textit{Backend}}

As operações documentadas para criação, edição, exclusão e consulta de publicações são relacionadas às histórias correspondentes na Tabela~\ref{tab:endpoints-feed}.

\renewcommand{\arraystretch}{1.5}
\setlength{\LTleft}{\fill}
\setlength{\LTright}{\fill}
\begin{longtable}{|>{\centering\arraybackslash}p{1.2cm}|>{\centering\arraybackslash}p{1.7cm}|>{\raggedright\arraybackslash}p{5.7cm}|>{\raggedright\arraybackslash}p{5.2cm}|}
\caption{\textit{Endpoints} do feed.}\label{tab:endpoints-feed}\\
\hline
\textbf{US} & \textbf{Método} & \textbf{Rota} & \textbf{Operação} \\
\hline
\endfirsthead
\multicolumn{4}{c}{\tablename\ \thetable{} -- continuação}\\
\hline
\textbf{US} & \textbf{Método} & \textbf{Rota} & \textbf{Operação} \\
\hline
\endhead
\hline
\endfoot
\hline
\multicolumn{4}{c}{\small Fonte: autoria própria.}\\
\endlastfoot
US21 & \textit{POST} & \texttt{/publicacoes} & Criar publicação. \\
\hline
US22 & \textit{PUT} & \texttt{/publicacoes/\allowbreak \{id\}} & Atualizar publicação. \\
\hline
US23 & \textit{DELETE} & \texttt{/publicacoes/\allowbreak \{id\}} & Excluir publicação. \\
\hline
US24 & \textit{GET} & \texttt{/publicacoes} & Listar publicações para visualização do \textit{feed}. \\
\hline
US24 & \textit{GET} & \texttt{/publicacoes/\allowbreak \{id\}} & Consultar publicação por identificador. \\
\hline
\end{longtable}
